% Generated by paper-covert from output/manuscript/sections_revised/08_limitations_and_future_work.md

\section{Limitations and Future Work}\label{limitations-and-future-work}

\subsection{Limitations, in order of severity}\label{limitations-in-order-of-severity}

\begin{enumerate}
\def\labelenumi{\arabic{enumi}.}
\tightlist
\item
  \textbf{Area over-call (1.47-1.59x), reproduced nationally and regionally}, is the map's most consequential operating characteristic (Section 5.2, Section 6.1). It is mitigated, but not eliminated, by shipping \texttt{ndvi\_amp} and \texttt{abstain\_reason} on every polygon (Section 6.4), which lets a downstream user filter toward a tighter, more area-accurate subset when that matters more than completeness.
\item
  \textbf{Legume misclassification} --- 18.4\% mapped versus 8.7\% ABS share, with the argmax-versus-mean-probability check ruling out a decision-rule artefact (Section 5.3, discussed as the paper's largest unresolved classification error in Section 6.7) --- has no established mechanism yet and remains untouched by this round of work.
\item
  \textbf{\texttt{unsegmented\_blob}}, SAM's single largest segmentation failure mode, accounts for 48,016 polygons nationally; a size-aware second segmentation pass to split these was never attempted (Section 4.1).
\item
  \textbf{Single-season national coverage.} The completed national map covers 2024 only; a national multi-year extension (2017-2023, 2025) is planned but not yet executed {[}PENDING: E8{]}. This is currently the paper's largest scope gap relative to its own framing as a multi-year resource, and the reason the multi-year results in Section 5.6 and Figure 7 are explicitly regional rather than national.
\item
  \textbf{No canola yield.} An optical-only model's coefficient of determination (0.27) did not exceed a trivial year-and-state-mean baseline (0.29); we judged shipping a yield estimate that underperforms guessing the season and region to be worse than shipping none, and disclose this rather than silently omitting canola yield without explanation.
\item
  \textbf{Sentinel-1's value is task-dependent, not a single verdict.} Re-tested against the current best classifier candidate (Section 6.5, Section 7.1) rather than the retired baseline it was originally measured against, Sentinel-1 regresses classification accuracy ($-$0.023 macro F1 temporal, concentrated on Legume) but clearly improves the shipped Cereal yield model (+0.092 R\textsuperscript{2} on its operational transfer bar; Table 7). Neither result changes what ships today --- the deployed classifier uses neither Sentinel-1 nor the candidate feature set, and no Sentinel-1 product is used at national scale for either task --- but the earlier, single-direction ``weak positive gain, not worth the cost'' framing of this limitation no longer describes the evidence and is corrected here.
\item
  \textbf{Three classes, not the full ten-species roster} originally scoped. A nine-species classifier reached only macro F1 0.38 at the available field-verified label volume; this is a finding about label-volume requirements for field-verified species classification (Section 1, Section 4.2), stated here for completeness and stated earlier, in the Introduction, so that a reader encounters it as a design decision rather than discovers it only in the Results.
\item
  \textbf{The ABS validation reference is itself uncertain, and is not an Olofsson-style probability-sample area estimate.} ABS and ABARES disagree with each other by a median 6.8\% on national sown area; no map-versus-ABS comparison in this paper should be read as more precise than that floor. No probability sample of field-verified crop-species labels exists for Australia at the density an Olofsson-style \citep{olofsson2014goodpractices} stratified area estimate with confidence intervals would require, so this paper's validation is a composition-and-area comparison against an independent statistic, not a stratified accuracy assessment (Section 2.4, Section 4.6).
\item
  \textbf{The WorldCereal and NLUM spatial comparisons (Section 5.5) carry a three-year reference-year mismatch} (2021 versus this map's 2024) and, for WorldCereal, cover only the presence and Cereal classes (Canola and Legume have no counterpart in WorldCereal's class scheme); both comparisons were run on a stratified sample rather than the full polygon population, adequate for the proportions and agreement rates reported but not for an area-total comparison.
\item
  \textbf{A leading classifier candidate exists that outperforms the shipped model, and has not been adopted.} An internally tested candidate feature set --- the shipped model's three indices plus six band-derived indices from \citet{sharma2026wa} --- reaches macro F1 0.890 (temporal) and 0.892 (spatial), against the shipped model's 0.821/0.823, with no change to the operating-point canola detection rate (89.0\% at 5\% false-positive rate; Table 8). Every classifier number reported elsewhere in this paper describes the shipped model, not this candidate. We disclose it here rather than only in project records because a reviewer with access to the accompanying code would reasonably ask why a better-performing configuration was not used: the reason is process, not evidence. This candidate has not yet been through an independent review separate from the process that produced it --- a check this project applies before adopting any model change --- and adopting it would additionally require re-running the national map generation and ABS validation this paper reports, both of which used the shipped model throughout (Section 7.2).
\end{enumerate}

\begin{longtable}[]{@{}
  >{\raggedright\arraybackslash}p{(\linewidth - 6\tabcolsep) * \real{0.2500}}
  >{\raggedright\arraybackslash}p{(\linewidth - 6\tabcolsep) * \real{0.2500}}
  >{\raggedright\arraybackslash}p{(\linewidth - 6\tabcolsep) * \real{0.2500}}
  >{\raggedright\arraybackslash}p{(\linewidth - 6\tabcolsep) * \real{0.2500}}@{}}
\toprule\noalign{}
\begin{minipage}[b]{\linewidth}\raggedright
Metric
\end{minipage} & \begin{minipage}[b]{\linewidth}\raggedright
Shipped (3 indices, 51 features)
\end{minipage} & \begin{minipage}[b]{\linewidth}\raggedright
Candidate (shipped + Sharma et al.~indices, 153 features)
\end{minipage} & \begin{minipage}[b]{\linewidth}\raggedright
Change
\end{minipage} \\
\midrule\noalign{}
\endhead
\bottomrule\noalign{}
\endlastfoot
Macro F1, temporal & 0.821 & 0.890 & +0.069 \\
Macro F1, spatial & 0.823 & 0.892 & +0.069 \\
Canola detected @ 5\% false-positive rate & 89.0\% & 89.0\% & +0.0 pp \\
Canola average precision & 0.944 & 0.935 & $-$0.009 \\
Canola ROC AUC & 0.965 & 0.962 & $-$0.003 \\
Independent review completed & Yes (this is the shipped model) & Not yet & --- \\
\end{longtable}


\textbf{Table 8.} Leading classifier candidate versus the shipped model. Reported here as a disclosed, not-yet-adopted finding; no result elsewhere in this paper uses the candidate.

\subsection{Future work, in priority order}\label{future-work-in-priority-order}

\begin{enumerate}
\def\labelenumi{\arabic{enumi}.}
\tightlist
\item
  \textbf{Complete the national 2017-2025 multi-year run (E8)} --- the direct next step, and the one this paper's own multi-year claims are explicitly written to accommodate once it lands. This step followed a discussion of the current validation methodology and the consensus/multi-year regional approach with project colleagues, using an earlier version of this manuscript as part of that discussion; the run is in progress at the time of this draft, with two of the nine planned seasons complete.
\item
  \textbf{Investigate the Legume over-prediction mechanism} --- the largest unexplained classification error in the map (Limitation 2).
\item
  \textbf{A size-aware second SAM segmentation pass} targeting the \texttt{unsegmented\_blob} failure mode (Limitation 3).
\item
  \textbf{Validate two-pass presence gating} (exempting Canola from the phenology-shape gate) against the ABS area ratio. This variant restores canola recall to 94.5\% at no measured cost to Cereal or Legume recall in the presence-recall check alone (Section 5.4), but has not yet been checked against the area-ratio metric that decided the original phenology-shape gate's rejection, and should not be treated as validated until it has.
\item
  \textbf{Extend the WorldCereal/NLUM spatial comparison from a 250,000-point sample to the full polygon population}, and to additional independent products, once the reference-year mismatch can be narrowed (e.g., against a future WorldCereal release closer to 2024) or a sensitivity check against reference year is added.
\item
  \textbf{Independently review the leading classifier candidate} (Limitation 10) and, contingent on that review and a project decision, either adopt it --- which would require re-running the national map generation and ABS validation reported in this paper --- or retain it as a documented, not-adopted alternative.
\item
  \textbf{Extend Sentinel-1 acquisition to national scale for the Cereal yield model specifically}, not the classifier (Section 6.5). This is gated on two unresolved items: whether the measured yield gain survives on ordinary commercial-paddock geometry rather than the clean National Variety Trial paddocks it was measured on --- no yield-labelled holdout on ordinary paddock geometry currently exists to check this --- and a revisit-density covariate for the Sentinel-1B multi-year acquisition discontinuity, which has not yet been implemented.
\end{enumerate}
